% Part: lambda-calculus
% Chapter: lambda-definability
% Section: fixpoints

\documentclass[../../../include/open-logic-section]{subfiles}

\begin{document}

\olfileid{lam}{ldf}{fp}

\olsection{స్థిరబిందువులు}

క్రమగుణిత ప్రమేయాన్ని $\fn{Fac}$ అనే పదంగా పునరావృత్తి ద్వారా,
కింది లక్షణంతో నిర్వచించాలనుకుందాం:
\[
\fn{Fac} \ident \lambd[n][\fn{IsZero}\, n\, \num 1 (\fn{Mult}\, n (\fn{Fac}(\fn{Pred} \, n)))]
\]
అంటే $n = 0$ అయితే $n$ క్రమగుణితం $1$; లేకపోతే అది
$n-1$ క్రమగుణితానికి $n$ రెట్లు. అయితే $\fn{Fac}$ను ఈ విధంగా
నిర్వచించలేం: కుడి వైపున $\fn{Fac}$ పదమే మళ్లీ వస్తుంది.
ఇలాంటి స్వీయ-సూచనతో కూడిన పునరావృత్త నిర్వచనాలు
లాంబ్డా కలనశాస్త్రపు పద నిర్వచనాలు కావు. ఉదాహరణకు ఒక పదాన్ని
\[
\fn{Mult} \ident \lambd[ab][a (\fn{Add}\, b) \num{0}]
\]
అని నిర్వచించినప్పుడు, కుడి వైపున ముందే నిర్వచించిన
$\fn{Add}$ వంటి పదాలే ఉంటాయి.
\sourcecorrection{OLTELAMLDFPIX-001}{ఆంగ్ల మూలపు గుణకార ఉదాహరణ $\lambd[ab][a(\fn{Add}\,a)0]$లో $b$ వాడబడదు; అది సాధారణ గుణకారం కాదు. ముందు అంకగణిత విభాగంలో ప్రకటించి సరిచేసిన ప్రత్యామ్నాయ గుణకార పదానికి అనుగుణంగా సంకలనంలో $b$ను, చర్చ్ శూన్యానికి $\num{0}$ను వాడాం.}
$\fn{Add}$ స్థానంలో దాని నిర్వచన పదాన్ని పెట్టి దాన్ని
తొలగించవచ్చు. అప్పుడు $\fn{Mult}$ శుద్ధ లాంబ్డా పదంగా
మారుతుంది. $\fn{Add}$లోనూ ఇతర నిర్వచిత పదాలు ఉంటే
(ఉదాహరణకు $\fn{Add}'$లో ఉన్నట్లు), ఇదే ప్రక్రియను
కొనసాగించి చివరికి శుద్ధ లాంబ్డా పదాన్ని పొందవచ్చు.

కానీ పై $\fn{Fac}$ వంటి పునరావృత్త నిర్వచనంలో అలా కాదు.
కుడి వైపున ఉన్న $\fn{Fac}$ స్థానంలో $\fn{Fac}$ నిర్వచనాన్నే
పెడితే ఇలా వస్తుంది:
\begin{align*}
  \fn{Fac} & \ident \lambd[n].\fn{IsZero}\, n\, \num{1} \\
    & \qquad (\fn{Mult}\, n
  ((\lambd[n][\fn{IsZero} \, n \, \num 1\, (\fn{Mult}\, n\, (\fn{Fac}
    (\fn{Pred}\, n)))]) (\fn{Pred}\, n)))
\end{align*}
\sourcecorrection{OLTELAMLDFPIX-002}{ఆంగ్ల మూలపు విస్తరణలో $\lambd[n][\fn{IsZero}\,n\,\num1]$ వద్ద లాంబ్డా పదం ముందుగానే ముగిసి, తరువాతి $\fn{Mult}$ శాఖ దాని వెలుపల పడింది. శూన్యం కాని సందర్భపు ఆ శాఖ కూడా అదే లాంబ్డా పదం శరీరంలో ఉండేలా $\lambd[n].$ రూపంతో రెండు పంక్తుల గణిత అమరికను పునరుద్ధరించాం; లోపలి $\fn{Fac}$ మాత్రం ఉద్దేశపూర్వకంగా మిగిలింది.}
ఇప్పటికీ కుడి వైపున $\fn{Fac}$ తొలగిపోలేదు. ఈ
ప్రక్రియను మళ్లీ చేస్తే నిర్వచనం ఇంకా పొడవవుతుంది;
ఎన్నిసార్లు చేసినా శుద్ధ లాంబ్డా పదం రాదు. కాబట్టి
క్రమగుణితాన్ని, సాధారణంగా పునరావృత్త ప్రమేయాలను,
ఈ పద్ధతిలో నిర్వచించడం సాధ్యం కాదు.

అయితే పునరావృత్త నిర్వచనం ఒక సూచన ఇస్తుంది. క్రమగుణిత
ప్రమేయాన్ని సూచించే పదం $f$ ఉంటే, ఈ పదాన్ని తీసుకుందాం:
\[
\fn{Fac}' \ident \lambd[g][\lambd[n][\fn{IsZero} \, n \, \num 1\, (\fn{Mult} \, n\, (g (\fn{Pred} n)))]]
\]
దాన్ని $f$కు ప్రయోగించిన $\fn{Fac}'\,f$ కూడా సంఖ్యలపై
క్రమగుణిత ప్రమేయాన్నే సూచిస్తుంది. అంటే $\fn{Fac}'$ ఒక
ప్రమేయాన్ని స్వీకరించి మరో ప్రమేయాన్ని ఇస్తుందని చూస్తే,
$f$ క్రమగుణితాన్ని సూచించే సందర్భంలో $\fn{Fac'}\, f$,
$f$ సంఖ్యాంక ఆర్గ్యుమెంట్లపై ఒకే ఫలితాలు ఇస్తాయి.
అంతమాత్రాన ఆ రెండు పదాలు బీటా-తుల్యాలని తేలదు.
\sourcecorrection{OLTELAMLDFPIX-003}{ఆంగ్ల మూలం ఒకే సంఖ్యాప్రమేయాన్ని సూచించడం నుంచి $\fn{Fac}'f$ పదం $f$కే సమానమని, ఆ వెంటనే బీటా-స్థిరబిందువు అని తేల్చింది. సంఖ్యాంకాలపై సమాన ఫలితాలు రావడం మాత్రమే పదాల $\beta$-తుల్యతను నిర్ధారించదు. కాబట్టి ప్రాతినిధ్యాన్ని, తరువాతి బలమైన స్థిరబిందు సమీకరణాన్ని వేరుగా చెప్పాం; ఆ సమీకరణకు పరిష్కారాన్ని కింది $Y$ నిర్మాణం ఇస్తుంది.}
$\fn{Fac}' \, f \equal[\beta] f$ అనే బలమైన లక్షణం
ఉన్న పదం~$f$ను $\fn{Fac}'$కు \emph{స్థిరబిందువు} అంటాం.
అలాంటి పదం క్రమగుణితానికి $\fn{Fac}'$ స్థిరబిందు ప్రాతినిధ్యం ఇస్తుంది.

లాంబ్డా కలనశాస్త్రంలో ఇచ్చిన పదానికి స్థిరబిందువులను
నిర్మించే పదాలు ఉన్నాయి. వాటి ద్వారా $\fn{Fac}'$ వంటి
పదాన్ని క్రమగుణిత నిర్వచనంగా మార్చవచ్చు.

\begin{defn}
  \ollabel{defn:Turing-Y}
  \emph{Y-సంయోజకం} ఈ పదం:
  \[
  Y \ident (\lambd[ux][x(uux)])(\lambd[ux][x(uux)]).
  \]
\end{defn}

\begin{thm}
  ఏ పదం $g$కైనా $Yg \red g(Yg)$ అనేది $Y$ లక్షణం.
  కాబట్టి $Yg$ ఎప్పుడూ~$g$కు స్థిరబిందువు.
\end{thm}

\begin{proof}
  $(\lambd[ux][x(uux)])$ను~$U$ అని సంక్షిప్తంగా రాద్దాం;
  అప్పుడు $Y \ident UU$. కనుక
  \begin{align*}
    Y g &\ident (\lambd[ux][x(uux)])U\,g \\
    &\red (\lambd[x][x(UUx)])g\\
    & \red g(UUg) \ident g(Yg).
  \end{align*}
  $g(Yg)$, $Yg$ రెండూ $g(Yg)$కే తగ్గుతాయి. అందువల్ల
  $g(Yg) \equal[\beta] Yg$; కాబట్టి $Yg$ అనేది~$g$కు
  స్థిరబిందువు.
\end{proof}

$Yg$లో రెడెక్స్ ఉన్నందున తగ్గింపును నిరవధికంగా
కొనసాగించవచ్చు:
\sourcecorrection{OLTELAMLDFPIX-004}{ఆంగ్ల మూలం $Yg$ స్వయంగా ఒక రెడెక్స్ అని చెప్పింది. $Y=UU$ విస్తరణలో $Yg=(UU)g$కి వెలుపలి ప్రయోగం రెడెక్స్ కాదు; లోపలి $UU$ రెడెక్స్. అందువల్ల పదంలో రెడెక్స్ ఉందని మాత్రమే సరిచేశాం; కింది తగ్గింపు క్రమం మారలేదు.}
\begin{align*}
  Y g &\red g (Y g) \\
    &\red g (g (Y g)) \\
    &\red g(g (g (Y g)))\\
    &\ldots
\end{align*}
అందువల్ల $Yg$ను $g$ అనంతసార్లు వరుసగా
ప్రయోగించబడిన రూపంగా ఊహించవచ్చు. దానికి~$g$ను
మరొకసారి ప్రయోగించినా, ఒక విధంగా చెప్పాలంటే,
అదనంగా ఏమీ చేయనట్లే: $Yg$లోని $g$ను
అనంతసార్లు ప్రయోగించిన~$Yg$కు మళ్లీ~$g$ను
ప్రయోగించినా, $g$ ప్రయోగాల అనంత పునరావర్తన రూపమే ఉంటుంది.

పై~$Yg$ నుంచి మొదలయ్యే $\beta$-తగ్గింపు క్రమం
అనంతమని గమనించండి. కాబట్టి $Yg$ను ఏదైనా పదానికి
ప్రయోగించి $(Yg)N$ను పరిశీలించినా, అది
$(g(Yg))N$, $(g(g(Yg)))N$, \dots\ వంటి అనంతమైన
వేర్వేరు పదాలకు తగ్గుతుంది. అయినా, \emph{మరో}
తగ్గింపు క్రమం నియత రూపంలో ఆగే అవకాశం ఉంది.

ఉదాహరణకు క్రమగుణితాన్ని తీసుకుందాం. $\fn{Fac}$ను
$Y\, \fn{Fac}'$గా, అంటే $\fn{Fac'}$కు స్థిరబిందువుగా,
నిర్వచిద్దాం. అప్పుడు:
\begin{align*}
  \fn{Fac}\,\num 3
  &\red Y\, \fn{Fac}' \, \num 3 \\
  &\red \fn{Fac}' (Y \,\fn{Fac}') \, \num 3 \\
  &\ident (\lambd[x][\lambd[n][\fn{IsZero} \, n\, \num 1\,
      (\fn{Mult}\, n\, (x (\fn{Pred}\, n)))]]) \, \fn{Fac} \, \num 3\\
  & \red \fn{IsZero} \, \num 3\, \num 1\,
      (\fn{Mult}\, \num 3\, (\fn{Fac} (\fn{Pred}\, \num 3))) \\
  &\red  \fn{Mult}\, \num 3 \, (\fn{Fac} \, \num 2).
  \intertext{అదేవిధంగా,}
  \fn{Fac}\,\num 2
  &\red  \fn{Mult}\, \num 2 \, (\fn{Fac} \, \num 1) \\
  \fn{Fac}\,\num 1
  &\red  \fn{Mult}\, \num 1 \, (\fn{Fac} \, \num 0)
  \intertext{కానీ}
  \fn{Fac}\,\num 0
  &\red \fn{Fac}' (Y \,\fn{Fac}') \, \num 0 \\
  &\ident (\lambd[x][\lambd[n][\fn{IsZero} \, n\, \num 1\,
      (\fn{Mult}\, n\, (x (\fn{Pred}\, n)))]]) \, \fn{Fac} \, \num 0\\
  & \red \fn{IsZero} \, \num 0\, \num 1\,
      (\fn{Mult}\, \num 0\, (\fn{Fac} (\fn{Pred}\, \num 0))).\\
  &\red \num 1.
  \intertext{కలిపి చూస్తే}
  \fn{Fac}\,\num 3
  &\red  \fn{Mult}\, \num 3 \,
  (\fn{Mult} \, \num 2\, (\fn{Mult} \, \num 1\, \num 1)).
\end{align*}

$\fn{Fac'}$కు వర్తించిన పద్ధతి ఏ పునరావృత్త నిర్వచనానికైనా
వర్తిస్తుంది. ఒక పునరావృత్త సమీకరణ ఉందనుకుందాం:
\begin{align*}
g\,x_1\dots x_n & \equal[\beta] N
\intertext{ఇక్కడ $N$లో~$g$, $x_1$, \dots,~$x_n$ ఉండవచ్చు. అప్పుడు
$G \ident (Y \lambd[g][\lambd[x_1\dots x_n][N]])$ అనే పదం ఎప్పుడూ
ఉంటుంది; దానికి}
G\,x_1 \dots x_n & \equal[\beta] \Subst{N}{G}{g}.
\intertext{ఎందుకంటే స్థిరబిందు సిద్ధాంతం ప్రకారం,}
G \ident (Y \lambd[g][\lambd[x_1\dots x_n][N]]) & \red \lambd[g][\lambd[x_1\dots x_n][N]](Y \lambd[g][\lambd[x_1\dots x_n][N]])\\
& \ident (\lambd[g][\lambd[x_1\dots x_n][N]])\,G
\intertext{కాబట్టి}
G\,x_1 \dots x_n
& \red (\lambd[g][\lambd[x_1\dots x_n][N]])\,G\,x_1\dots x_n\\
& \red (\lambd[x_1\dots x_n][\Subst{N}{G}{g}])\,x_1\dots x_n\\
  & \red \Subst{N}{G}{g}.
\end{align*}

\olref{defn:Turing-Y}లోని $Y$-సంయోజకాన్ని అలన్ ట్యూరింగ్
ప్రతిపాదించాడు. అలోంజో చర్చ్ మరో రూపాన్ని ప్రతిపాదించాడు;
దాన్ని~$Y_C$ అంటాం:
\[
Y_C \ident \lambd[g][(\lambd[x][g(xx)])(\lambd[x][g(xx)])].
\]
చర్చ్ సంయోజకం ట్యూరింగ్ సంయోజకం కంటే కొంత బలహీనమైనది:
$Y_Cg \equal[\beta] g(Y_Cg)$ అయినా,
$Y_Cg \bred g(Y_Cg)$ కాదు.
\sourcecorrection{OLTELAMLDFPIX-005}{ఆంగ్ల మూలం చర్చ్ $Y_C$ను ట్యూరింగ్ $Y$తో పోలుస్తూ రెండు సమీకరణాల్లోనూ $Yg$ అని వ్రాసింది. కానీ ముందరి సిద్ధాంతం $Yg\red g(Yg)$ని ఇప్పటికే నిరూపించింది; ఇక్కడ బీటా-తుల్యత మాత్రమే ఉన్నది $Y_Cg$కు. కాబట్టి ఆ రెండు సమీకరణాల్లో $Y_C$ను పునరుద్ధరించాం.}
$V$ అనేది $\lambd[x][g(xx)]$ పదం అనుకుందాం;
అప్పుడు $Y_C \ident \lambd[g][VV]$. కనుక
\begin{align*}
  VV & \ident (\lambd[x][g(xx)])V \red g(VV)
  \text{ కనుక}\\
  Y_C g & \ident (\lambd[g][VV])g \red VV \red g(VV), \text{ అయితే}\\
  g(Y_C g) & \ident g((\lambd[g][VV])g) \red g(VV).
\end{align*}
మరో మాటలో, $Y_Cg$, $g(Y_Cg)$ రెండూ ఒకే
$g(VV)$ పదానికి తగ్గుతాయి. కనుక
$Y_Cg \equal[\beta] g(Y_Cg)$. అనేక అనువర్తనాలకు
ఇది సరిపోతుంది.

\end{document}
